\section{六、异质性分析}\label{ux516dux5f02ux8d28ux6027ux5206ux6790}

为考察数据要素利用的定价效率改善效应在不同制度环境下的差异，本文在统一滞后规格下报告三条仍具有较清晰解释力的异质性结果：政策阶段，以及高分析师关注和上市板块两条补充性边界条件。所有模型均控制企业与年份固定效应和11个控制变量，标准误在行业×年份水平聚类。政策阶段结果同时结合交互项检验和费舍尔组合检验，高分析师关注与上市板块则以交互项检验为主。

\subsection{（一）政策阶段效应}\label{ux4e00ux653fux7b56ux9636ux6bb5ux6548ux5e94}

相较于2020或2021，2022年``数据二十条''更接近数据要素制度密集推进的起点，也更适合作为主文中的政策阶段切点。其背后的逻辑并非机械地将估计结果解释为某一单一政策的即时效应，而是将2022年之后理解为``数据利用披露逐步制度化、市场预期更充分''的阶段。若数据利用信息在较早阶段更具新颖性和辨识度，则其对价格发现的边际贡献应在2022年前更强。

表12第(1)(2)列结果显示，2011至2021年样本中DU\_kw系数为-0.0066（t=-5.13），2022至2024年样本中系数为-0.0018（t=-0.85，不显著），交互项检验P值为0.031，费舍尔组合检验得到相同方向的结论。这意味着数据要素利用降低股价延迟的效应主要集中在制度化推进之前的阶段：在相关披露尚未成为更普遍表达时，年报中的数据利用信息更容易构成有效的信息更新；而当市场对该类叙事已有更高先验预期后，其边际定价价值有所减弱。

\textbf{表12 政策阶段异质性检验}

{\def\LTcaptype{} % do not increment counter
\begin{longtable}[]{@{}lll@{}}
\toprule\noalign{}
& (1) & (2) \\
\midrule\noalign{}
\endhead
\bottomrule\noalign{}
\endlastfoot
变量 & 2011-2021 & 2022-2024 \\
& PriceDelay & PriceDelay \\
DU\_kw & -0.0066*** & -0.0018 \\
& (0.0013) & (0.0021) \\
控制变量 & YES & YES \\
Firm FE & YES & YES \\
Year FE & YES & YES \\
N & 24,720 & 12,240 \\
R² & 0.455 & 0.557 \\
交互项P值 & 0.031** & \\
\end{longtable}
}

替代切点检验进一步表明，以2021为切点或剔除2020过渡年后再比较前后阶段，组间差异均不稳定，因此本文不再将2020或2021作为主文中的首选断点。相关结果见附录。

\subsection{（二）信息中介活跃度}\label{ux4e8cux4fe1ux606fux4e2dux4ecbux6d3bux8dc3ux5ea6}

除制度阶段外，信息中介的活跃程度也可能影响数据利用披露的吸收速度。若分析师跟踪更充分，年报中的相关表述更容易被解释、传播并进入价格，数据利用披露的边际定价价值因而可能更高。

表13第(1)(2)列显示，高分析师关注组中，DU\_kw的系数为-0.0048（t=-4.18），低分析师关注组中为-0.0025（t=-1.75，仅边际显著），交互项P值为0.001。该结果提供了一条补充性边界证据：信息中介越活跃，年报中的数据利用信息越容易被资本市场转化为更快的价格发现。这一证据与ForecastDisp渠道形成互补：前者强调分析师群体对同一企业信息解释更趋一致，后者说明在分析师本身更活跃的企业中，相关披露的定价含量也更容易被释放。

\subsection{（三）上市板块差异}\label{ux4e09ux4e0aux5e02ux677fux5757ux5deeux5f02}

上市板块差异则反映了企业所处信息环境和投资者预期的不同。主板公司经营模式相对成熟、投资者覆盖更广，年报中的数据利用披露未必形成强增量信号；相较之下，非主板企业通常成长性更强、业务不确定性更高，数据利用信息更可能成为投资者重新理解企业价值的重要线索。

表13第(3)(4)列显示，主板公司中DU\_kw的系数为-0.0017（t=-0.98），并不显著；非主板公司中系数为-0.0042（t=-4.01），在1\%水平上显著为负，交互项P值小于0.001。该结果同样提供补充性边界证据：数据要素利用披露的边际定价价值在非主板企业中更强，说明当信息摩擦和成长性不确定性更高时，年报中的相关表述更容易构成有效的信息更新。

\textbf{表13 补充异质性检验}

{\def\LTcaptype{} % do not increment counter
\begin{longtable}[]{@{}lllll@{}}
\toprule\noalign{}
& (1) & (2) & (3) & (4) \\
\midrule\noalign{}
\endhead
\bottomrule\noalign{}
\endlastfoot
维度 & 高分析师关注 & 低分析师关注 & 主板公司 & 非主板公司 \\
因变量 & PriceDelay & PriceDelay & PriceDelay & PriceDelay \\
DU\_kw & -0.0048*** & -0.0025* & -0.0017 & -0.0042*** \\
& (0.0011) & (0.0014) & (0.0017) & (0.0010) \\
控制变量 & YES & YES & YES & YES \\
Firm FE & YES & YES & YES & YES \\
Year FE & YES & YES & YES & YES \\
N & 18,912 & 17,847 & 15,900 & 21,394 \\
交互项P值 & 0.001*** & & \textless0.001*** & \\
\end{longtable}
}

在统一滞后规格下，数字经济核心产业口径的交互项P值为0.507，未能提供稳定的组间差异证据；因此，本文不再将其作为主文异质性的正式结论，仅在附录中作为补充结果简要报告。

\subsection{（四）行业披露水平的同群溢出效应}\label{ux56dbux884cux4e1aux62abux9732ux6c34ux5e73ux7684ux540cux7fa4ux6ea2ux51faux6548ux5e94}

前三条异质性证据说明数据要素利用披露的定价效应受制度阶段、信息中介活跃度和上市板块的调节。一个互补的问题是：除企业自身披露之外，所处行业的整体披露水平是否独立地影响企业股价的信息吸收速度。若行业整体披露水平更高，同业之间的可比性和信息解读规范可能更成熟，投资者在评估单个企业时所面临的解释成本也相应下降，从而形成``同群溢出''。

为此，本文构造行业-年份层面的披露渗透率指标IndPen\_mean，定义为企业所在二级行业在当年的DU\_kw均值，并采用留一法剔除企业自身。将该指标与企业自身DU\_kw同时放入主回归，可分别识别``自身披露''与``行业信息环境''两条渠道的效应。为使系数在量纲上可比，两项指标均作标准化处理。控制变量、固定效应结构和聚类方式与主文一致。

表14第(1)列显示，单独纳入企业自身DU\_kw\_z时，系数为-0.0086（t=-5.10），与主文基准一致。第(2)列同时纳入IndPen\_mean\_z后，企业自身DU\_kw\_z仍显著为负（-0.0067，t=-5.26），而行业披露渗透率IndPen\_mean\_z的系数为-0.0087（t=-2.47），在5\%水平上显著。第(3)列进一步在滞后规格下重复该回归，结果显示行业渗透率的系数增大至-0.0158（t=-4.46），在1\%水平上高度显著，绝对值甚至超过企业自身披露的系数（-0.0070，t=-3.55）。这表明，行业整体披露水平与企业股价延迟之间存在独立且稳健的负向关联，且在时间顺序更清晰的滞后规格下更为突出。

\textbf{表14 行业披露水平的同群溢出效应}

{\def\LTcaptype{} % do not increment counter
\begin{longtable}[]{@{}llll@{}}
\toprule\noalign{}
& (1) & (2) & (3) \\
\midrule\noalign{}
\endhead
\bottomrule\noalign{}
\endlastfoot
规格 & 基准 & 同期 & 滞后 \\
因变量 & PriceDelay & PriceDelay & PriceDelay \\
DU\_kw\_z & -0.0086*** & -0.0067*** & \\
& (0.0017) & (0.0013) & \\
DU\_kw\_z\_lag & & & -0.0070*** \\
& & & (0.0020) \\
IndPen\_mean\_z & & -0.0087** & \\
& & (0.0035) & \\
IndPen\_mean\_z\_lag & & & -0.0158*** \\
& & & (0.0035) \\
控制变量 & YES & YES & YES \\
Firm FE & YES & YES & YES \\
Year FE & YES & YES & YES \\
N & 43,312 & 43,312 & 38,061 \\
\end{longtable}
}

该结果为主文``政策阶段效应''提供了补充性理解。2022年之后，数据要素披露的个体边际效应减弱，但行业整体披露水平仍与更快的价格发现相关。两条证据共同表明：单家企业披露的边际新颖性随制度化推进而下降，但行业层面披露规范的成熟带来了更稳定的信息环境改善。据此，``数据要素利用披露的定价价值''不应被理解为个体层面的单一信号效应，而应理解为``个体披露的相对新颖性''与``行业披露形成的信息共同体''共同作用的结果。需要说明的是，进一步引入企业自身DU\_kw与行业渗透率的交互项时，其系数在两种渗透率度量下方向不一致，本文因此不将其作为主文结论，相关结果见附录。

总体而言，主文异质性证据首先支持``信息新颖性与市场先验预期''这一主要边界条件，并进一步提供``信息中介活跃度''``上市板块信息摩擦''和``行业披露水平的同群溢出''三条补充性边界证据，而非广泛的产业差异叙事。
